
(Volume II. The Congregation.)

The second volume deals with the congregation. It contains 6 sections, or 25 pieces, 9 of which are unsigned editorial articles from Sverdrup’s pen. For local reasons, almost all these treatises and contributions were composed in a popular form; the presentation is always clear and transparent; but it lies in the nature of the matter that not a little requires labor on the reader’s part. Here too the rule applies: “The one who has, to him shall be given.” The one who comes to the book with questions will receive answers.

It is not our intention to enter into the individual articles. We shall mention most of them and attach remarks to such as seem to us to have a particular significance for the understanding of the ecclesiastical development among us.


First Section.

1. “What Is at Stake?”, written as early as 1877, with reference to the controversy carried on between the Synod and the Conference, shows that even a whole generation before Sverdrup laid down his work, it was the congregation that was at issue: “Our definite conviction is this, that in this controversy it is the congregation that is at stake.”

2. “The Holiness of the Congregation” states clearly that the holiness of the congregation, and the question bound up with it, “What is congregation?”, constitutes “the true chief question in the Free Church.” It should be noted, first, that the article was written in 1889, the year before the union, and many years before the Lutheran Free Church was formed; second, that Sverdrup was opposed to the “unsound views” for which Pastor Trandberg and his Reformed compatriots were spokesmen.

It is emphasized that there may be congregation without organization and constitution; for the congregation in Jerusalem was not organized; further, that the empirical congregation is congregation, even though many of its members are hypocrites, dead, sleeping, unbelieving, and ungodly. Sverdrup’s position here—as it always was—is in agreement with Article 7 of the Augsburg Confession, which depicts the congregation according to its idea, true essence, its invisible-visible form—and with Article 8, which depicts the congregation according to its historical-empirical appearance, a mixture of hypocrites and evil men with the believers, who constitute the proper substance of the congregation.

3. “The Liberation of the Congregation,” 1896, points out that it was the new direction in the Conference that “held back,” so that “the congregation’s disempowerment” did not proceed so rapidly, and that the old direction found staunch allies in the union, so that the United Church in the course of a few years left “all the other bodies far behind in denominational domination and tyranny over the congregations.” The author lays the blame for this partly upon the State Church upbringing. “Many times that which in the State Church was done by compulsion is here among us imitated by habit and custom. Popular custom is a stronger power than the law for compelling men, and the old habits reign long after legal compulsion has ceased, especially when many zealous pastors labor with all their might to maintain the old habit and custom, that all shall belong to the congregation” (p. 19). Among other things, S. emphasizes the necessity of lay activity and points out that in the Norwegian State Church there is quite great freedom for lay activity because the law protects it, and that there was far greater freedom for lay activity there than in most “majority congregations.”

Rank and class distinctions, official grandeur, must go.

4. “The Vivification of the Congregation,” 1898, contains, among other things, not a little on Luther’s view of polity.

Second Section.

1. “Which Church Polity Is Most Fitting for the Evangelical Lutheran Church?”, 1869. This is perhaps the oldest literary contribution from Sverdrup’s hand that we now possess. It is interesting because it gives us an insight into what a theological student, 20 years old, could say about church polity. One will be astonished that Sverdrup at so young an age could embrace an ecclesiastical constitutional question with such deep interest as this treatise bears witness to, and still more at the ability with which it is handled. As far as I can see, S. is here the spokesman for Rudelbach’s ideas, in that he pleads for consistories and synods, but disapproves of independent congregations.

In our country Sverdrup was, from first to last, the defender of the independent congregations. For the free land he found that the episcopal, the presbyterial, the consistorial polity that one had in the Lutheran lands of Europe stood far behind the congregational. The Free Church was therefore built upon a congregationalist basis. That was something new—at any rate in the history of the Lutheran Church. In Europe too eyes have been opened to the superiority of a free-congregation polity. None less than the clear-sighted and learned church historian Loofs at the University of Halle says: “Who knows whether the congregationalist church polity does not also have its future among us, when one day the State Churches in the old world collapse.”

2. “The Free Church Association,” 1877–78. This treatise, of 46 pages, I regard as the most important in the whole book. It falls into six subdivisions: The Congregation, The Local Congregation, The Office in the Free Congregation, The Pastoral Office and Lay Activity, The Diaconal Office in the Congregation, The Fellowship of the Congregations. Here we receive answers to questions such as: how shall we establish congregations? what does it mean “to compel them to come in”?

It is the Word, says S.—not outward compulsion or means of enticement or party spirit—that shall gather a congregation. “Even among the Norwegian emigrants, who indeed had all once been received into God’s Church through Baptism, it would have been far better if one had never organized any congregation or gathered members of a congregation in any other way or by any other means than by the preaching of God’s Word, publicly and privately. For even baptized persons are many a time, alas, in such a condition that it is better both for themselves and for the congregation that they ‘stand outside’ and are the object of the congregation’s influence in that way, than that they should be inside and perhaps there sleep the deepest sleep, precisely because they are ‘in the congregation.’”

After dwelling somewhat upon the all too hasty admissions into the congregation, and the dangers bound up with church discipline and exclusion, Sverdrup lays the weight upon this: that the congregation must be gathered by God’s Word, and that the constraining and drawing and calling of the Word is the only thing that is to be used in order to bring people into the congregation.

One will ask: “Should a man be admitted into the congregation if he is not converted?” To this S. replies: “If God’s Spirit is permitted to work this resolve in a man, that he desires admission into the congregation, then he is a real increase of the congregation, even if he has not come to the peace of faith with God; without this he is only a diminishment of the congregation’s true strength.”

“There are even within the little circle where we are acquainted congregations whose existence we fear is due to party spirit and hatred and purely carnal motives. And there are many members of the congregation entered in the rolls who would never have been there, if openly outward, purely carnal advantages had not been used as bait to get them in. There has not everywhere been the proper respect for the people’s freedom, nor either the proper reverence for the congregation’s purity. It is therefore natural that such congregations will reap what they have sown.”

Very important, viewed from the dogmatic-ethical standpoint as well as from the practical-theological, are also the statements about “The Office in the Free Congregation” and “The Pastoral Office and Lay Activity.” Here an account is given of the difference between the pastoral office and the universal spiritual priesthood. How much confusion would not our congregations be spared if they made themselves familiar with what Sverdrup writes about the office. We are afraid that there are also among the pastors those who are not clear on the fact that the office is instituted by God, but not the clerical estate. How often can a pastor not become a tyrannizing hireling of the congregation, or the dictator of an immature congregation.

There is, however, one small section in which I do not agree with Sverdrup. It is concerning the origin of the diaconal office. He takes the point of departure for this in Acts, chapter 6. The congregation in Jerusalem is supposed to have chosen deacons. Alongside the deacons the presbyters are supposed to have labored. The other apostolic congregations are supposed to have arranged themselves after the congregation in Jerusalem.


The seven chosen men (Acts 6:3) are never called deacons in the New Testament. Philip (Acts 21:8) is called simply one of the seven. Their work is not called diakonia either. The task of the seven men was different from that which the deacons in the other congregations had. The work of the deacons was to care for the poor; the seven merely determined who should receive support; the deacons made no such determination. The requirements set for the seven (Acts 6:3) are stricter than those made of the deacons (1 Tim. 3:6). Stephen, one of the seven, was chiefly a preacher; and Philip, another of the seven, was an evangelist. The seven stood chiefly in relation to the apostles, and their work had much in common with that which we later learn was performed by the presbyters. Much speaks for the view that the seven were presbyters.

The first time deacons are mentioned in the New Testament is in the Epistle to the Philippians 1:1 (about the year 63) and 1 Tim. 3:8 (a couple of years later). The new translation calls them servants of the congregation.

That the congregation in Jerusalem was a pattern for the congregations in Judea is likely. That its polity was imitated in the Pauline congregations is not demonstrable. Neither the synagogue, nor the temple, nor the Greek associations determined the polity, but local conditions and practical demands. Then as now polity was a matter of freedom.

3. “Is It Possible for Norwegian Lutherans to Build a Free Christian Church in America?”, 1880–81, is a reckoning with some theologians in Norway who said that the Conference reminded them all too often that it belonged in the Land of Humbug. Sverdrup draws out the consequence of their statements: the spirit that presses for the freedom of the congregation is the mark of humbug. He does not leave the answer owing.

4. “Free Congregation in a Free Church,” 1882, is an answer to the declaration of the Thirty, which was submitted by 30 members at the Conference’s annual meeting in 1882. Here we learn from what is cited from the declaration of the Thirty that it was “The Open Declaration,” composed in 1874 by Professors Oftedal and Weenaas, that was the point of departure for the new direction.

5. “Denominational Power and the Free Congregation,” 1897–98, is directed against “the rule of the majority” in the United Church. Pastor J. A. Kildahl receives considerable attention in this piece. Among other things, it is stated that when Augsburg split, and the United Church transferred three professors away, it let the other seven remain behind without even asking whether they wished to move with it.
